% GENERATED FILE. Edit canonical lessons, then run npm run generate:books.
% canonical-source-hash: fnv1a64:34084a8eaffb6907
% canonical-lessons: MW-C107-bunno, MW-C107-rangno, MW-C107-lagavno, MW-C107-hatavno, MW-C107-bichhavno

\chapter{To weave, To dye, To apply, To remove, To spread out}
\label{ch:mw-to-weave-to-dye-to-apply-to-remove-to-spread-out}
\hlchaptermodality{107}

\begin{chapteropening}
\textbf{By the end of this chapter:} \emph{I can say to weave, to dye, to apply, to remove and to spread out.}

Complete the last lesson of chapter 107: I can say to weave, to dye, to apply, to remove and to spread out.
\end{chapteropening}

\section[\texorpdfstring{\mw{बुणणो} (buṇṇo)}{बुणणो (buṇṇo)}]{\mw{बुणणो} (buṇṇo) — to weave}

\label{lesson:MW-C107-bunno}



\begin{quote}
Before the new one: say the Marwadi for to burn, then the Marwadi for to sew.
\end{quote}

\subsection*{You'll want to know: \mw{बुणणो}}
\textbf{\mw{बुणणो}} — \emph{buṇṇo} — \textquotedblleft{}to weave\textquotedblright{}.

\begin{grammarlens}[title={What you've built}]
One more word for this chapter.
\end{grammarlens}

\subsection*{Your turn}
\begin{itemize}
\raggedright
  \item \emph{Say it:} \emph{buṇṇo}
  \item \emph{Say it:} \emph{buṇṇo}, once more
  \item \emph{Say it:} say \emph{buṇṇo}
  \item \emph{Recall:} say the Marwadi for to burn, then the Marwadi for to sew, then say \emph{buṇṇo} again
\end{itemize}

\begin{tcolorbox}[breakable,colback=teal!4,colframe=teal!35!black,title={Before you move on}]
What does \mw{बुणणो} mean? (\textquotedblleft{}To weave\textquotedblright{}.) Say it once more.
\end{tcolorbox}
\section[\texorpdfstring{\mw{रंगणो} (raṅgṇo)}{रंगणो (raṅgṇo)}]{\mw{रंगणो} (raṅgṇo) — to dye}

\label{lesson:MW-C107-rangno}



\begin{quote}
Before the new one: say the Marwadi for to sew, then the Marwadi for to weave.
\end{quote}

\subsection*{You'll want to know: \mw{रंगणो}}
\textbf{\mw{रंगणो}} — \emph{raṅgṇo} — \textquotedblleft{}to dye\textquotedblright{}.

\begin{grammarlens}[title={What you've built}]
One more word for this chapter.
\end{grammarlens}

\subsection*{Your turn}
\begin{itemize}
\raggedright
  \item \emph{Say it:} \emph{raṅgṇo}
  \item \emph{Say it:} \emph{raṅgṇo}, once more
  \item \emph{Say it:} say \emph{raṅgṇo}
  \item \emph{Recall:} say the Marwadi for to sew, then the Marwadi for to weave, then say \emph{raṅgṇo} again
\end{itemize}

\begin{tcolorbox}[breakable,colback=teal!4,colframe=teal!35!black,title={Before you move on}]
What does \mw{रंगणो} mean? (\textquotedblleft{}To dye\textquotedblright{}.) Say it once more.
\end{tcolorbox}
\section[\texorpdfstring{\mw{लगावणो} (lagāvṇo)}{लगावणो (lagāvṇo)}]{\mw{लगावणो} (lagāvṇo) — to apply, to fix on}

\label{lesson:MW-C107-lagavno}



\begin{quote}
Before the new one: say the Marwadi for to weave, then the Marwadi for to dye.
\end{quote}

\subsection*{You'll want to know: \mw{लगावणो}}
\textbf{\mw{लगावणो}} — \emph{lagāvṇo} — \textquotedblleft{}to apply, to fix on\textquotedblright{}.

\begin{grammarlens}[title={What you've built}]
One more word for this chapter.
\end{grammarlens}

\subsection*{Your turn}
\begin{itemize}
\raggedright
  \item \emph{Say it:} \emph{lagāvṇo}
  \item \emph{Say it:} \emph{lagāvṇo}, once more
  \item \emph{Say it:} say \emph{lagāvṇo}
  \item \emph{Recall:} say the Marwadi for to weave, then the Marwadi for to dye, then say \emph{lagāvṇo} again
\end{itemize}

\begin{tcolorbox}[breakable,colback=teal!4,colframe=teal!35!black,title={Before you move on}]
What does \mw{लगावणो} mean? (\textquotedblleft{}To apply, to fix on\textquotedblright{}.) Say it once more.
\end{tcolorbox}
\section[\texorpdfstring{\mw{हटावणो} (haṭāvṇo)}{हटावणो (haṭāvṇo)}]{\mw{हटावणो} (haṭāvṇo) — to remove}

\label{lesson:MW-C107-hatavno}



\begin{quote}
Before the new one: say the Marwadi for to dye, then the Marwadi for to apply.
\end{quote}

\subsection*{You'll want to know: \mw{हटावणो}}
\textbf{\mw{हटावणो}} — \emph{haṭāvṇo} — \textquotedblleft{}to remove\textquotedblright{}.

\begin{grammarlens}[title={What you've built}]
One more word for this chapter.
\end{grammarlens}

\subsection*{Your turn}
\begin{itemize}
\raggedright
  \item \emph{Say it:} \emph{haṭāvṇo}
  \item \emph{Say it:} \emph{haṭāvṇo}, once more
  \item \emph{Say it:} say \emph{haṭāvṇo}
  \item \emph{Recall:} say the Marwadi for to dye, then the Marwadi for to apply, then say \emph{haṭāvṇo} again
\end{itemize}

\begin{tcolorbox}[breakable,colback=teal!4,colframe=teal!35!black,title={Before you move on}]
What does \mw{हटावणो} mean? (\textquotedblleft{}To remove\textquotedblright{}.) Say it once more.
\end{tcolorbox}
\section[\texorpdfstring{\mw{बिछावणो} (bichhāvṇo)}{बिछावणो (bichhāvṇo)}]{\mw{बिछावणो} (bichhāvṇo) — to spread out}

\label{lesson:MW-C107-bichhavno}



\begin{quote}
Before the new one: say the Marwadi for to apply, then the Marwadi for to remove.
\end{quote}

\subsection*{You'll want to know: \mw{बिछावणो}}
\textbf{\mw{बिछावणो}} — \emph{bichhāvṇo} — \textquotedblleft{}to spread out\textquotedblright{}.

\begin{grammarlens}[title={What you've built}]
One more word for this chapter.
\end{grammarlens}

\subsection*{Your turn}
\begin{itemize}
\raggedright
  \item \emph{Say it:} \emph{bichhāvṇo}
  \item \emph{Say it:} \emph{bichhāvṇo}, once more
  \item \emph{Say it:} say \emph{bichhāvṇo}
  \item \emph{Recall:} say the Marwadi for to apply, then the Marwadi for to remove, then say \emph{bichhāvṇo} again
\end{itemize}

\begin{tcolorbox}[breakable,colback=teal!4,colframe=teal!35!black,title={Before you move on}]
What does \mw{बिछावणो} mean? (\textquotedblleft{}To spread out\textquotedblright{}.) Say it once more.
\end{tcolorbox}
